\subsection{Actual rigid discs and a designated low hole}\label{subsec:fc-line-affinoid-punctures}

\begin{lemma}[Closed discs form a basis at rigid points]\label{lem:fc-rigid-closed-disc-basis}
\lean{EnvelopeProperty.Berkovich.RelativeSpectrum.fc_exists_closedDisc_rigid_nhds}
\leanok
\uses{thm:fc-separated-affinoid-basis, lem:fc-inflation-linear}
For every actual rigid evaluation point $a$ and every neighbourhood $N$
of it, there is $r>0$ such that the whole actual closed disc $D(a,r)$
is contained in $N$.
\end{lemma}

\begin{proof}
Take an actual affinoid with the rigid point in its strict core, contained
in $N$. Choose $r$ smaller than the outer radius and all its positive
distances from the hole centres. The ultrametric linear-coordinate
identities show every point of $D(a,r)$ satisfies every affinoid inequality.
\end{proof}

\begin{lemma}[Adding an actual separated hole]\label{lem:fc-separated-affinoid-puncture}
\lean{EnvelopeProperty.Berkovich.RelativeSpectrum.FCStandardAffinoid.puncture, EnvelopeProperty.Berkovich.RelativeSpectrum.FCStandardAffinoid.IsSeparated.puncture, EnvelopeProperty.Berkovich.RelativeSpectrum.FCStandardAffinoid.puncture_carrier, EnvelopeProperty.Berkovich.RelativeSpectrum.FCStandardAffinoid.puncture_strictCore, EnvelopeProperty.Berkovich.RelativeSpectrum.FCStandardAffinoid.puncture_boundary}
\leanok
\uses{def:fc-separated-affinoid-core}
Let $U$ have a separated presentation, let the rigid point $a$ lie in its
strict core, and let $r>0$ be smaller than the outer radius and all distances
from $a$ to the old hole centres. Removing the actual open disc $D^-(a,r)$
gives an actual separated standard affinoid $V$. Its carrier and strict core
are respectively those of $U$ intersected with $\{|T-a|\ge r\}$ and
$\{|T-a|>r\}$. Its actual boundary is the union of the old boundary and
the actual point $\zeta_{a,r}$.
\end{lemma}

\begin{proof}
Append the original centre $a$ and radius $r$ to the actual hole data.
The old hole radii are smaller than their distances from $a$ because the
rigid point lies in the strict core. The new radius is also smaller than
these distances by construction. The ultrametric centre-distance bound
therefore proves the new open hole disjoint from every old open hole,
and all strict separation inequalities hold for the same presentation.
The set identities follow directly from these literal appended data.
\end{proof}
